%############## ТАБЛИЦА ##############%
\begin{table}[H]
    \caption{Полученные экспериментальные и расчетные величины.}
    % \label{tab1}
    \centering
    \begin{tabular}{|c|cc|ccc|c|}
        \hline 
    \end{tabular}
\end{table}

%############## КАРТИНКА ##############%
\begin{figure}[H]
    \centering
    \includegraphics[width=0.7\textwidth, frame]{plots/norm.eps}
    \caption{Оценка нормальной ширины щели.}
    % \label{fig2}
\end{figure}

%############## МНОГО КАРТИНОК ##############%
\begin{figure}[!htb]
    \centering
    
    \subfloat[Диафрагма min]{ \includegraphics[frame, clip,width=0.33\columnwidth]{plots/sp_min.eps} }
    \subfloat[Диафрагма 10 мм]{ \includegraphics[frame, clip,width=0.33\columnwidth]{plots/sp_10.eps} }
    
    \subfloat[Диафрагма 30 мм]{ \includegraphics[frame, clip,width=0.33\columnwidth]{plots/sp_30.eps} }
    
    \caption{Зарегистрированные спектры при разных диафрагмах.}
    % \label{fig3}
\end{figure}

%############## УРАВНЕНИЕ ##############%
\begin{flalign} 
    % \label{eq1}
    a
\end{flalign}

%############## ПОДСТРАНИЦЫ ##############%
\begin{figure}[H]
    \begin{minipage}[!t]{0.49\linewidth}
        
    \end{minipage}    
    \hfill
    \begin{minipage}[!t]{0.49\linewidth}
        
    \end{minipage}    
\end{figure}